\section{The binary model and the proof framework}\label{sec:framework}

We first reduce the singularity probability to the expected size of a set
of possible new columns. We then introduce the potential that will control
this set and state the two estimates on which the rest of the proof rests.

\subsection{From signs to bits}

A \emph{binary matrix} of size \(m\) is a symmetric \(m\times m\) matrix
whose entries on and above the diagonal are independent and uniform on
\(\bits\). We write \(\mathbf B_m\) for a random binary matrix and use
ordinary letters for its realizations.

\begin{lemma}\label{lem:sign-bit}
For \(n\ge2\),
\[
 \Pr(\det A_n=0)=\Pr(\det\mathbf B_{n-1}=0).
\]
\end{lemma}

\begin{proof}
Put \(\epsilon=(A_n)_{11}\), multiply \(A_n\) by \(\epsilon\), and then
apply the diagonal congruence with diagonal
\[
 1,\ \epsilon(A_n)_{12},\ldots,\epsilon(A_n)_{1n}.
\]
The resulting first row and column consist of ones. Conditional on the
original first row, the lower principal block still has independent
uniform signs on and above its diagonal: every one of these signs has
only been multiplied by a sign determined by the first row. Its Schur
complement at the top left entry is therefore \(-2\mathbf B_{n-1}\).
The two operations and the Schur complement preserve singularity.
\end{proof}

For the remainder of the proof, \(N\) is a reference integer tending to
infinity and
\[
 \ell=\log\log N,\qquad
 \mathcal I_N=\{m\in\Z:N-N^{7/10}\le m\le N\}.
\]
The logarithm \(\log\) is natural, while \(\log_2\) and \(\log_N\) have
the indicated bases. Every constant exponent of \(\ell\) is fixed before
\(N\) tends to infinity. In particular,
\(\ell^D/\log N=o(\ell^{-A})\) for all fixed \(A,D>0\).
We write \([m]=\{1,\ldots,m\}\).

We say that a property of \(\mathbf B_m\) \emph{holds typically} if its
failure probability is at most \(2^{-Nf(N)}\), uniformly for
\(m\in\mathcal I_N\), for some \(f(N)\to\infty\). When the property is
stated only for nonsingular matrices, failure means nonsingularity
together with failure of the stated conclusion. The typical set may
depend on any fixed precision parameters in the statement. A finite
intersection of typical properties is again typical.

All unqualified ranks are over \(\R\). We will also use the following
fact over finite fields.

\begin{lemma}[Principal rank]\label{lem:principal-rank}
Let \(M\) be a symmetric matrix over a field, of rank \(r\). Then \(M\)
has a nonsingular principal submatrix of size \(r\).
\end{lemma}

\begin{proof}
Choose a matrix \(X\) whose \(r\) columns form a basis of the column
space of \(M\), and choose a left inverse \(T\) of \(X\). Since \(M\) is
symmetric, its row space is the row space of \(X^{\mathsf T}\). Thus
\[
 M=X S X^{\mathsf T},\qquad S=TMT^{\mathsf T}.
\]
The matrix \(S\) is symmetric and has rank \(r\). Choose \(r\) independent
rows of \(X\). On those indices the principal submatrix of \(M\) is a
product of three nonsingular \(r\times r\) matrices.
\end{proof}

\begin{lemma}[Corank tail]\label{lem:corank-tail}
Let \(\mathbb K\) be any field, and regard the bits of \(\mathbf B_m\)
as elements of \(\mathbb K\). For \(1\le k\le m\),
\[
 \Pr\bigl(\corank_{\mathbb K}\mathbf B_m\ge k\bigr)
 \le \sum_{j=k}^{m}\binom mj\,2^{-j(j+1)/2}.
\]
Consequently, uniformly for \(m\in\mathcal I_N\),
\[
 \Pr\bigl(\corank_{\mathbb K}\mathbf B_m>N^{3/5}\bigr)
 \le 2^{-\Omega(N^{6/5})}.
\]
The implicit constant in the last bound is independent of the field.
\end{lemma}

\begin{proof}
If the rank is \(m-j\), Lemma~\ref{lem:principal-rank} supplies a
nonsingular principal submatrix on some \(m-j\) indices. Reveal all
entries in the corresponding columns. The remaining symmetric
\(j\times j\) block is forced over \(\mathbb K\), since its Schur
complement must vanish. Each forced entry has at most one possible bit
value, so the conditional probability is at most \(2^{-j(j+1)/2}\).
Union over the choices of indices and over \(j\ge k\). For
\(j>N^{3/5}\), the factor \(2^{-j(j+1)/2}\) dominates the at most
\(2^m\) choices of indices and the \(m\) choices of \(j\).
\end{proof}

\begin{lemma}[Reduction to corank one]\label{lem:corank-reduction}
For the binary model in dimension \(N\),
\[
 \Pr(\det\mathbf B_N=0)
 \le 2^{o(N)}\Pr(\corank\mathbf B_N=1)
       +2^{-\Omega(N^{6/5})}.
\]
\end{lemma}

\begin{proof}
Let \(B\) have corank \(k\ge1\), and write \(K=\ker B\). Choose \(k-1\)
coordinates whose restrictions to \(K\) are independent linear
functionals. Flip the diagonal bit at each of those coordinates, and
write \(D\) for the diagonal matrix of the changes. Its selected
diagonal entries are all nonzero.

If \((B+D)v=0\), pairing this equation with every \(u\in K\) gives
\(u^{\mathsf T}Dv=0\). The chosen coordinate functionals on \(K\) are
independent, so the selected coordinates of \(v\) are zero. It follows
that \(Bv=0\). Conversely, every vector in \(K\) with these coordinates
zero lies in \(\ker(B+D)\). This subspace has dimension one. Hence
flipping \(k-1\) diagonal bits produces a matrix of corank one.

For every matrix with \(k\le N^{3/5}\), fix one such choice of
coordinates. Each resulting matrix has at most
\[
 \sum_{j\le \lceil N^{3/5}\rceil}\binom Nj
   =2^{O(N^{3/5}\log N)}=2^{o(N)}
\]
preimages, because a preimage is recovered by specifying the flipped
coordinates. All binary matrices are equally likely. The matrices of
larger corank have the probability bounded in
Lemma~\ref{lem:corank-tail}.
\end{proof}

\subsection{Admissible columns}

For a nonsingular binary matrix \(B\) of size \(m\), define
\begin{equation}\label{eq:admissible-set}
 \mathcal S(B)=
 \left\{z\in\bits^m:
 z^{\mathsf T}B^{-1}z\in\bits,\quad
 \|B^{-1}z\|_\infty\le1\right\}.
\end{equation}
Set \(\mathcal S(B)=\varnothing\) when \(B\) is singular. These are the
columns that can occur when a corank-one matrix is reduced by deleting
a largest coordinate of its kernel vector.

\begin{lemma}\label{lem:column-reduction}
One has
\[
 \Pr(\corank\mathbf B_N=1)
 \le N\,2^{-N}\,\E|\mathcal S(\mathbf B_{N-1})|.
\]
\end{lemma}

\begin{proof}
Suppose a symmetric matrix \(M\) of size \(N\) has corank one, and
choose \(0\ne v\in\ker M\). Delete an index \(i\) for which \(|v_i|\)
is largest. The resulting principal submatrix \(B\) is nonsingular.
Indeed, the adjugate of \(M\) is a nonzero scalar multiple of
\(vv^{\mathsf T}\), so its \(i\)-th diagonal entry is nonzero.

Write \(z\) for the column at \(i\) on the remaining indices and \(a\)
for its diagonal bit. The kernel equation and the Schur complement give
\[
 B^{-1}z=-v_{[N]\setminus\{i\}}/v_i,\qquad
 a=z^{\mathsf T}B^{-1}z.
\]
Thus \(z\in\mathcal S(B)\). For any fixed deleted index, conditional on
\(B\), the pair \((z,a)\) is uniform on \(\bits^{N-1}\times\bits\).
For each \(z\in\mathcal S(B)\), the displayed equality fixes \(a\).
The probability for that index is therefore
\(2^{-N}\E|\mathcal S(\mathbf B_{N-1})|\). Union over \(i\).
\end{proof}

It remains to prove
\begin{equation}\label{eq:column-target}
 \E|\mathcal S(\mathbf B_{N-1})|\le 2^{o(N)}.
\end{equation}
Indeed, Lemmas~\ref{lem:corank-reduction} and
\ref{lem:column-reduction} then give the upper bound \(2^{-N+o(N)}\)
for the binary singularity probability. A specified row of
\(\mathbf B_N\) is zero with probability \(2^{-N}\), which gives the
matching lower bound.

\subsection{A potential on the discriminant group}

For nonsingular \(B\), the finite abelian group
\[
 G_B=\Z^m/B\Z^m
\]
has order \(|\det B|\). It records the possible denominators of
\(B^{-1}z\) for integer \(z\). We now give a numerical measure of its
invariant factors.

For a finite abelian group \(G\), write its \(p\)-primary part as
\[
 G_p\cong\bigoplus_{i\ge1}\Z/p^{b_{p,i}}\Z,\qquad
 b_{p,1}\ge b_{p,2}\ge\cdots\ge0,
\]
where only finitely many exponents are positive. For \(e\ge1\), put
\[
 k_G(p,e)=\#\{i:b_{p,i}\ge e\},\qquad
 \nu_p=\frac{\log_N p}{N}.
\]
We call \((p,e)\) with \(k_G(p,e)>0\) a \emph{layer}, its integer
\(k_G(p,e)\) its \emph{height}, and \(\nu_p\) its \emph{weight}.
All sums over layers below are finite. The identity
\begin{equation}\label{eq:group-order}
 \sum_{p,e} k_G(p,e)\nu_p=\frac{\log|G|}{N\log N}
\end{equation}
follows from \(\sum_e k_G(p,e)=\sum_i b_{p,i}\).
Hadamard's inequality gives \(|\det B|\le N^{N/2}\) for a binary
matrix of size \(m\le N\). Hence the sum in
Equation~\eqref{eq:group-order} is at most \(1/2\) for \(G_B\) and
for every quotient of \(G_B\).

Write \(k_B(p,e)=k_{G_B}(p,e)\), set \(\beta=7/5\), and define
\begin{equation}\label{eq:potential}
 \Psi(B)=\sum_{p,e}\nu_p\,\phi(k_B(p,e)),\qquad
 \phi(0)=\phi(1)=0,\quad
 \phi(2)=\beta,\quad
 \phi(k)=2k\quad(k\ge3).
\end{equation}
In particular \(0\le\Psi(B)\le1\), since \(0\le\phi(k)\le2k\).
We set \(\Psi(B)=0\) for singular \(B\) when an expression also
contains a nonsingularity indicator.

For example, a \(p\)-primary part with positive cyclic exponents
\((5,3,1)\) has
\[
 (k_B(p,1),\ldots,k_B(p,5))=(3,2,2,1,1),
\]
so this primary part contributes \((6+2\beta)\nu_p\) to \(\Psi(B)\).
Figure~\ref{fig:layers} depicts these layers. The coefficient
\(\beta=7/5\) exceeds the coefficient \(5/4\) in the height-two term of
the color-cost bound in Section~\ref{sec:discriminant}, while satisfying
the scalar budget inequality in Section~\ref{sec:potential-step}.

The potential is chosen to support two estimates. The first relates
the number of admissible columns to the arithmetic structure of \(B\).

\begin{proposition}[Admissible-column estimate]\label{prop:admissible-potential}
For every fixed \(A>0\), typically for nonsingular binary matrices
\(B\) of size \(m\in\mathcal I_N\),
\[
 |\mathcal S(B)|\le 2^{N(\Psi(B)+\ell^{-A})}.
\]
\end{proposition}

The second estimate controls the potential in expectation.

\begin{proposition}[Potential moment]\label{prop:path-potential}
For a binary matrix of size \(N-1\),
\[
 \E\!\left[\ind_{\{\det\mathbf B_{N-1}\ne0\}}\,
            2^{N\Psi(\mathbf B_{N-1})}\right]\le 2^{o(N)}.
\]
\end{proposition}

These propositions imply Equation~\eqref{eq:column-target}. Indeed,
fix \(A>0\) in Proposition~\ref{prop:admissible-potential}. On its
typical set use that estimate, and on the exceptional set use
\(|\mathcal S(B)|\le2^N\). For a suitable \(f(N)\to\infty\),
\[
 \E|\mathcal S(\mathbf B_{N-1})|
 \le 2^{N\ell^{-A}}
     \E\!\left[\ind_{\{\det\mathbf B_{N-1}\ne0\}}\,
                2^{N\Psi(\mathbf B_{N-1})}\right]
       +2^{N-Nf(N)}
 =2^{o(N)}.
\]

Sections~\ref{sec:lattice-tools}--\ref{sec:discriminant} prove the static
estimate in Proposition~\ref{prop:admissible-potential} by combining the
discriminant pairing with lattice, rank, and spectral estimates.
Sections~\ref{sec:minor-moves}--\ref{sec:potential-step} prove a local
conditional estimate for extensions: certificate pricing handles width
one, while Proposition~\ref{prop:admissible-potential} directly pays for
the gate at width two. Section~\ref{sec:minor-paths} adds a separate
high-layer tail argument and combines these estimates along paths to
prove Proposition~\ref{prop:path-potential}.
